% Originally: content/week-12.tex, \section{Orientable manifolds} (Corsin Week 12)

\section{Orientable manifolds}
\label{sec:orientable_manifolds}

% Source: Corsin Nick/Class Notes/Week 12.pdf, p. 7
% Corsin reproduces the two definitions below directly from the lecture script, rather than
% writing them out in his own hand, and adds margin annotations (tutor reference moved out of the
% typeset text, 2026-11-04).
\begin{ainote}
The two definitions below are those of the lecture script (Definitions 14.68 and 14.65).
\end{ainote}

\begin{definition}[Support of a $k$-form]
\label{def:support_of_k_form}
% Correction: Claude Opus 5 (Ultracode) --- dropped a second German mirroring of the term
% "Traeger" from this line. The mirroring rule is first-introduction-only, and
% def:support_of_function in the preceding section already carries it; a German term appearing
% twice is a style error rather than a harmless repetition, per the maintenance note at the head
% of appendix-b-glossary.tex. (Spelled without the macro here on purpose, so that a grep for the
% mirroring macro does not report this comment as a third occurrence.)
Let $U \subseteq \mathbb{R}^n$ be open and $\alpha \in \Omega^k(U)$. The \newterm{support}
of $\alpha$ is the closed set
\[\supp\alpha := \overline{\{x\in U : \alpha_x \neq 0\}} \subseteq \overline{U}.\]
We say $\alpha$ has \newterm{compact support} in $U$ if $\supp\alpha$ is bounded and contained
in $U$.
\end{definition}

\begin{definition}[Orientable manifold]
\label{def:orientable_manifold}
Let $M \subseteq \mathbb{R}^n$ be a $k$-dimensional $C^\infty$ submanifold. We say $M$ is
\newterm{orientable} \germanfor{orientierbar} if there exists a family of local
parametrizations $\phi_i : V_i \to M \cap U_i$, $i \in I$, whose images cover $M$, and such
that for every $i,j \in I$ with $M \cap U_i \cap U_j \neq \emptyset$, the transition map
$\widetilde\phi_i^{-1}\circ\widetilde\phi_j : \widetilde V_j \to \widetilde V_i$ has positive
Jacobian determinant:
\[\det D(\widetilde\phi_i^{-1}\circ\widetilde\phi_j) > 0 \qquad \text{on } \widetilde V_j.\]
Here the tilde denotes restriction to the overlap: $\widetilde V_j := \phi_j^{-1}(M\cap U_i\cap
U_j)$, $\widetilde\phi_j := \phi_j|_{\widetilde V_j}$, and likewise for $i$.
Such a family $\{\phi_i\}$ is called an \newterm{atlas} of $M$; if it satisfies the
determinant condition above it is called an \newterm{orientation} of $M$. Once such a family
is chosen, $M$ is said to be \newterm{oriented}.
\end{definition}

% Moved out of the typeset text (the book names tutors only in source citations):
% Corsin marks the family {phi_i} itself with the margin note "is called an Atlas" next to the
% covering condition, and the positivity condition with "Atlas" again next to the determinant
% inequality, flagging that "atlas" names the family regardless of whether the determinant
% condition holds, while "orientation" is the stronger notion once it does. This matches the
% definition's own wording above. The figure below redraws his hand-drawn counterexample.

The figure below shows how the orientation condition fails:
two overlapping local parametrizations $\phi, \psi$ of the same patch of a surface, whose
coordinate frames $(\partial_1\phi,\partial_2\phi)$ and $(\partial_1\psi,\partial_2\psi)$ at a
shared point disagree in handedness.

% Sonnet 5 (Medium) --- FIG-W12-01
% Revised 2026-08-09 (Opus 5, high). One correction and one clarification.
%
%   * THE CORRECTION. The figure drew a rotation where it needed a reflection, i.e. it
%     illustrated the opposite of what the ainote below it asserts. With the frames based at
%     (2.1,1.3), the old vectors were
%         d1 phi = (1.2, 0.25),  d2 phi = (-0.4, 0.85):  cross = 1.02 + 0.10 = +1.12
%         d1 psi = (-1.2,-0.25), d2 psi = (0.45,-0.95):  cross = 1.14 + 0.1125 = +1.2525
%     Both cross products are POSITIVE, so both frames turned counterclockwise and the psi
%     frame was just the phi frame rotated by 180 degrees --- a transition map of positive
%     determinant, which is precisely the case an orientable atlas allows. d2 psi now points
%     at 55 degrees instead, giving d2 psi = (0.55, 0.78) and
%         cross = (-1.2)(0.78) - (-0.25)(0.55) = -0.936 + 0.1375 = -0.7985 < 0,
%     so the psi frame turns clockwise, as the ainote says it does. 55 degrees was chosen to
%     keep all four arrows apart: the directions are now 12, 115, 192 and 55 degrees, and the
%     tip (2.65,2.08) still lies inside the patch.
%   * THE CLARIFICATION. The two small frames on the right are drawn identically, which is
%     correct and was previously unexplained: in each parameter domain the frame is the
%     standard basis by definition, and the handedness mismatch appears only in the images.
%     They are now labelled as the domains, so the repetition reads as the point rather than
%     as an oversight.
% Feedback: Gemini 3.1 Pro (2026-08-11) --- called exceptional: two overlapping parametrizations
% whose frames disagree in handedness, tied straight to the non-orientability of the Möbius
% strip. Keep the frames disagreeing; the earlier version, a pure 180-degree rotation, drew the
% case an orientable atlas ALLOWS, and that regression is what the arithmetic above guards.
\begin{center}
\begin{tikzpicture}[>=Latex, scale=1.0,
  lbl/.style={fill=white, fill opacity=0.8, text opacity=1, inner sep=1.2pt}]
  \draw[green!50!black, thick] (0,0) .. controls (1.5,0.9) and (3.5,0.6) .. (5,1.2)
    .. controls (5,2.4) and (3,2.6) .. (1.5,2.0) .. controls (0,1.6) and (0,0.8) .. (0,0);
  \draw[green!50!black, dotted] (0.4,1.0) .. controls (1.8,0.6) and (3.2,1.1) .. (4.4,1.6);
  \fill (2.1,1.3) circle (1.2pt);
  \draw[blue!70!black, thick, ->] (2.1,1.3) -- (3.3,1.55) node[right, font=\scriptsize]
    {$\partial_1\phi$};
  \draw[blue!70!black, thick, ->] (2.1,1.3) -- (1.7,2.15) node[above, font=\scriptsize]
    {$\partial_2\phi$};
  % Correction: Claude Opus 5 (High) --- this label, `left' of the arrow tip at (0.9,1.05),
  % occupied roughly x = 0.15 to 0.9 at y = 1.05, and the dotted coordinate line starts at
  % (0.4,1.0) and runs just under y = 1.0 across exactly that span: the green dots printed
  % between the subscript and the psi. The free space left of the tip is a thin wedge between
  % the dotted line and the patch's own left edge (x = 0.14 at this height), too narrow to
  % move into, so the label takes the white plate from style.md instead.
  \draw[red, thick, ->] (2.1,1.3) -- (0.9,1.05)
    node[lbl, text=red, left, font=\scriptsize] {$\partial_1\psi$};
  \draw[red, thick, ->] (2.1,1.3) -- (2.65,2.08) node[right, font=\scriptsize] {$\partial_2\psi$};
  % Correction: Claude Opus 5 (High) --- these two named the parametrizations in plain black
  % while their frames are blue and red, so nothing tied phi to the blue frame or psi to the
  % red one; the reader had to infer it from position. Coloured to match, which is the same
  % pairing the two reference frames on the right already use.
  \node[blue!70!black, font=\scriptsize] at (0.1,-0.35) {$\phi$};
  \node[red, font=\scriptsize] at (4.4,2.7) {$\psi$};
  \draw[blue!70!black, ->] (6.3,0.2) -- (7.3,0.2) node[right, font=\scriptsize] {$e_1$};
  \draw[blue!70!black, ->] (6.3,0.2) -- (6.3,1.2) node[above, font=\scriptsize] {$e_2$};
  \node[blue!70!black, font=\scriptsize] at (6.8,-0.3) {domain of $\phi$};
  \draw[red, ->] (8.6,0.2) -- (9.6,0.2) node[right, font=\scriptsize] {$e_1$};
  \draw[red, ->] (8.6,0.2) -- (8.6,1.2) node[above, font=\scriptsize] {$e_2$};
  \node[red, font=\scriptsize] at (9.1,-0.3) {domain of $\psi$};
\end{tikzpicture}
\end{center}

% Corsin's caption on this figure: "Here, det(D(phi^{-1} o psi)) < 0!". The ainote that quoted it
% is now the remark below, without the tutor reference.
\begin{remark}[An atlas that is not an orientation]
At the shared point, $(\partial_1\phi,\partial_2\phi)$ turns counterclockwise while
$(\partial_1\psi,\partial_2\psi)$ turns clockwise, so the two frames are mirror images of each
other, and $\det(D(\phi^{-1}\circ\psi)) < 0$ there. This is exactly the atlas $\{\phi,\psi\}$
failing \cref{def:orientable_manifold}: it covers the patch, but its transition map has
\emph{negative} Jacobian determinant there, so $\{\phi,\psi\}$ is an atlas but not an
orientation.
\end{remark}

% Generator: Claude Sonnet 5 (Medium)
\begin{ainote}
\Cref{ex:stokes_work_along_paths} and Exercise 12.6 (the latter not reproduced below, see the priority table)
invoke \newterm{Stokes' theorem} in its classical $\mathbb{R}^3$ form, which this book derives
only in \cref{sec:stokes_in_r3}, as a consequence of the general theorem. It is recorded here
ahead of that, since the exercise-sheet solutions rely on it.
\end{ainote}

% Merged: Claude Opus 5.5 (Extra) --- this was stated a second time in
% content/26-stokes/04-stokes-theorem.tex as thm:stokes_3d (Paschoud's formulation, with the
% orientation convention spelled out). This statement is kept because ch. 23's solutions cite it,
% two chapters before ch. 26; it now carries the convention from thm:stokes_3d, whose citations
% point here. The derivation from thm:stokes_general stays in sec:stokes_in_r3.
\begin{theorem}[Stokes' theorem, classical form in $\mathbb{R}^3$]
\label{thm:stokes_classical}
% Generator: Claude Sonnet 5 (Medium)
% Source (orientation convention): Jérôme Paschoud/Notizen/Woche 13.1 Stokes.pdf, p. 4
Let $S \subseteq \mathbb{R}^3$ be a compact, oriented $C^1$ surface with boundary, oriented by a
continuous unit normal field $\nu$, and let $F$ be a $C^1$ vector field on a neighbourhood of
$S$. Let $\partial S$ carry the induced orientation, given by the unit tangent $\tau$ for which,
walking along $\partial S$ in the direction of $\tau$ with $\nu$ pointing up, $S$ lies to the
left (the \qt{right-hand rule}). Then
\[\int_S \langle\curl F,\nu\rangle\,d\vol_2 = \int_{\partial S} \langle F,\tau\rangle\,d\vol_1,
  \qquad\text{in physics notation}\qquad
  \int_S \curl \vec F \cdot d\vec A = \int_{\partial S} \vec F \cdot d\vec s.\]
\end{theorem}
